\begin{Huge}
    Bioinformatik-FAQ
\end{Huge}
\begin{block}{H\"aufig gestellte Fragen zum Studium}

\question{Lernt man im Studium, wie man programmiert?}{
    Ja, aber auf einer sehr eigenst\"andigen Basis. Man bekommt einen \"Uberblick \"uber die Sprache(n) und \"uber dahinterstehende Programmierkonzepte. Praktische Programmiererfahrung muss man sich aber gr\"o"stenteils selbst erarbeiten.
}

\question{Welche Programmiersprachen macht man da so?}{
    Das ist von Prof zu Prof und von Veranstaltung zu Veranstaltung unterschiedlich. In den ersten beiden Semestern lernt man meistens Racket und Java.
}

\question{Muss man programmieren k\"onnen, um das Studium anzufangen?}{
    Nein. Die Vorlesung beginnt absolut bei 0, um allen den Einstieg zu erm\"oglichen.
}

\question{Muss man gut in Mathe sein?}{
    Man muss kein Mathe-Genie sein, man sollte Mathe aber nicht hassen. Es ist gerade am Anfang viel Mathe.
}

\question{Muss ich Bio 4-st\"undig gehabt haben, um Bioinformatik zu studieren?}
{
    Nein, Bioinformatik hat mit der klassischen Schulbiologie wenig zu tun. Es ist allerdings hilfreich, wenn du schon mal das Schema einer Eukaryotenzelle gesehen hast (und wei"st, was Eukaryoten sind).
}

\question{Stehe ich als Bioinformatiker*in viel im Labor?}{
    Jein. Das Studium beinhaltet einige Laborpraktika, aber nicht ann\"ahernd so viel wie z.B. bei der Chemie oder Biochemie.
}

\question{Was ist der Unterschied zwischen Bio- und Medizininformatik?}{
    Die Bioinformatik besch\"aftigt sich grob gesagt mit auotmatisierter Verarbeitung von DNA, Molek\"ulstrukturen etc., Medizininformatik geht mehr in Richtung Patientendaten und medizinische Bildverarbeitung.
}

\question{Wie ist die Frauenquote so?}{
    44,3\%.
}

\question{Gibt es Praktika?}{
    Im normalen Studienverlauf ist kein berufsorientiertes Praktikum vorgesehen, viele arbeiten aber parallel als Werksstudent*in oder man macht ein Kurzpraktikum in den Semesterferien.
}

\question{Kann man ein Auslandssemster machen?}{
     Klar, geht immer. T\"ubingen nimmt am ERASMUS-Programm teil, die Organisation ist aber langwierig und man sollte sich fr\"uh (ein Jahr vorher) drum k\"ummern.
}

\question{Wie ist da so der NC?}{
    Die letzten 10 Jahre gab es keinen NC. Das muss aber theoretisch nicht immer so sein.
}
\end{block}
